\chapter*{Syllabus}
\addcontentsline{toc}{chapter}{Syllabus}

\hfill Academic year 2025/2026

\begin{center}\Large Gravitation\end{center}

\begin{enumerate}[label=\arabic*.,leftmargin=1.5em]
\item \textbf{Newtonian gravitation.}
  \begin{enumerate}[label=1.\arabic*.,leftmargin=2.5em,itemsep=0pt]
  \item Newton's law of universal gravitation.
  \item Limitations of Newtonian gravitation.
  \item Tidal forces.
  \end{enumerate}
\item \textbf{The equivalence principle.}
  \begin{enumerate}[label=2.\arabic*.,leftmargin=2.5em,itemsep=0pt]
  \item Gravitational mass and inertial mass. Active and passive gravitational mass.
  \item Gravitational forces and non-inertial forces.
  \item Galileo's equivalence principle.
  \item Einstein's equivalence principle: gravity and space-time geometry.
  \end{enumerate}
\item \textbf{Brief introduction to differential geometry.}
  \begin{enumerate}[label=3.\arabic*.,leftmargin=2.5em,itemsep=0pt]
  \item Differentiable manifold. Tensors and tensor fields. Forms.
  \item Lie derivative.
  \item Covariant derivative. Connections. Parallel transport.
  \item Curvature, torsion and non-metricity.
  \item Metric. Levi-Civita connection.
  \item Introductory tutorial to the xAct package for Mathematica.
  \end{enumerate}
\item \textbf{Matter in gravitational fields.}
  \begin{enumerate}[label=4.\arabic*.,leftmargin=2.5em,itemsep=0pt]
  \item Particle in a gravitational field.
  \item Locally inertial frames. Riemann normal coordinates.
  \item Geodesics. Proper time. Geodesic deviation equation.
  \item Symmetries. Killing vectors. Conserved quantities.
  \item Principle of general covariance. Electromagnetism in gravitational fields.
  \end{enumerate}
\item \textbf{Einstein's equations.}
  \begin{enumerate}[label=5.\arabic*.,leftmargin=2.5em,itemsep=0pt]
  \item Derivation of Einstein's equations. Einstein--Hilbert action.
  \item Gravity as an interaction mediated by a massless spin-2 particle. Coupling to matter. The equivalence principle as a consistency condition.
  \item Vacuum solutions. Maximally symmetric spaces.
  \item Cosmological solutions: Friedmann--Lemaître--Robertson--Walker metrics.
  \item Spherically symmetric solutions: Schwarzschild metric and black holes.
  \item Orbits in Schwarzschild: null and timelike geodesics. Classical tests of Einstein's equations.
  \end{enumerate}
\item \textbf{Gravitational waves.}
  \begin{enumerate}[label=6.\arabic*.,leftmargin=2.5em,itemsep=0pt]
  \item Linearisation of Einstein's equations.
  \item Gauge invariance.
  \item Emission of gravitational waves: the quadrupole radiation formula.
  \item Detection of gravitational waves: interferometers.
  \end{enumerate}
\end{enumerate}

\section*{Bibliography}
\begin{itemize}[label={},leftmargin=0pt,itemsep=0.6em]
\item Misner, Thorne and Wheeler, ``Gravitation'', Freeman, 1973.\\
This book contains a vast amount of information and is extremely useful as a reference book, but it is perhaps not the best option for a first approach to the subject.
\item Sean Carroll, ``Spacetime and Geometry: An Introduction to General Relativity'', Addison Wesley, 2004.\\
This book is shorter and more accessible, while still being complete. It is a good reference for a first approach. There is a set of lecture notes by Eugene Lim based on this book (E.~Lim was a PhD student of Sean Carroll).
\item James Hartle, ``Gravity: An Introduction to Einstein's General Relativity'', Pearson, 2002.\\
This book makes extensive use of physical intuition and introduces many of the concepts and developments from a more physical point of view. It is a highly recommended book.
\item Bernard Schutz, ``A First Course in General Relativity'', Cambridge University Press, 2009.\\
Another good book for a first approach to General Relativity. It starts from basic concepts of special relativity and then makes the step to general relativity in a very geometric but instructive way.
\item L.D.~Landau and E.M.~Lifshitz, ``The Classical Theory of Fields'' (Spanish edition: \emph{Teoría clásica de los campos}, second edition, Editorial Reverté, Barcelona, 1993).\\
As always, this is a book that is best read once one already has some familiarity with General Relativity. It also has a good number of solved exercises that are very illustrative and useful.
\item Robert Wald, ``General Relativity'', University of Chicago Press, Chicago, 1984.\\
Another of the great classics for studying General Relativity. A very formal and very rigorous book.
\item N.~Straumann, ``General Relativity. With Applications to Astrophysics'', Texts and Monographs in Physics, Springer, 2004.\\
This book is also very formal and requires being comfortable with a more advanced mathematical language.
\item Steven Weinberg, ``Gravitation and Cosmology: Principles and Applications of the General Theory of Relativity'', Wiley \& Sons, 1972.\\
This book approaches gravity in a way that differs from the more usual geometric treatments. Weinberg was one of the pioneers of quantum field theory and of the Standard Model of particle physics, and this shows in the exposition. It is highly recommended to study this book because it provides a very physical and complementary view.
\item Feynman Lectures on Gravitation, Addison-Wesley Publishing Company.\\
These lectures are also advisable reading, because Feynman was one of the pioneers of the formulation of General Relativity from a field-theory perspective. They may not be so useful for studying gravitation, but they are for seeing Feynman's perspective, which is always very enriching. For a more complete treatment, Weinberg is more advisable.
\item Michele Maggiore, ``Gravitational Waves'', Vol.~I.\\
This book is a treatise on the production, propagation and detection of gravitational waves. The first chapters treat the foundations of General Relativity in considerable detail. It has a good number of exercises and examples on gravitational waves that are very useful.
\item Ray D'Inverno, ``Introducing Einstein's Relativity'', Oxford University Press, 1992.
\item Eric Poisson, ``A Relativist's Toolkit: The Mathematics of Black-Hole Mechanics'', Cambridge University Press, 2004.
\item Eric Gourgoulhon, ``Special Relativity in General Frames'', Graduate Texts in Physics, Springer, 2013.
\end{itemize}
A couple of references for the differential geometry part:
\begin{itemize}[label={},leftmargin=0pt,itemsep=0.3em]
\item B.F.~Schutz, ``Geometrical Methods of Mathematical Physics'', Cambridge University Press, 1980.
\item Nakahara, ``Geometry, Topology and Physics'', Institute of Physics Publishing, IoP, 2003.
\item Y.~Choquet-Bruhat, C.~DeWitt-Morette, M.~Dillard-Bleick, ``Analysis, Manifolds and Physics'', Elsevier Science Publishing Company Inc., 1982.
\end{itemize}

In addition to the textbooks mentioned, a wide variety of lecture notes by professors who have taught General Relativity can be consulted. Among them, one can mention:
\begin{itemize}[itemsep=0.3em]
\item The lecture notes of Matthias Blau. An impressive compendium of General Relativity of more than 900 pages, where almost everything can be found with very clear explanations. It ranges from elementary concepts to very advanced topics, so it is suitable for all audiences.
\item The notes of Daniel Baumann are very didactic and easy to follow. Moreover, they are fairly short ($\sim100$ pages).
\item Bert Janssen, of the University of Granada, has written these General Relativity notes, which can serve as a very good guide for the course because their table of contents is quite similar to the syllabus.
\item These notes by Luis Garay, of the Complutense University of Madrid, are very good and contain very interesting material. However, they are more useful for seeing some extensions of what is covered in class than for following the course syllabus.
\end{itemize}
